\chapter{Higher generalizations}\label{ch:higher}

This chapter treats three generalizations: iterated cats, whose legs are
themselves cats; gaited cats, whose legs form a diagram rather than a family;
and two-dimensional cats in a $2$-category. In each case the analogue of
\cref{prop:automatic} holds, and the theory reduces to a known one.

\section{Iterated cats}\label{sec:iterated}

\begin{assumption}\label{ass:iterated}
In this section $\cC$ has an initial object $0$ and a terminal object $1$.
\end{assumption}

By \cref{prop:legless-initial,prop:terminal-cat}, $\Fel(\cC)$ has the initial
object $\Legless$ and the terminal object $\Caudal$, which we take as its chosen
extremal objects. Hence $\Fel(\Fel(\cC))$ is defined, and inductively we set
$\Fel^0(\cC)=\cC$ and $\Fel^{n+1}(\cC)=\Fel(\Fel^n(\cC))$. An object of
$\Fel^2(\cC)$ is a \emph{cat of cats}: a family of $\cC$-cats, equipped with a
legless head and a terminal tail.

\begin{proposition}\label{prop:flattening}
The category $\Fel(\cC)$ has small coproducts, so its body functor
\[
\body_{\Fel(\cC)}\colon\Fel^2(\cC)\too\Fel(\cC),\qquad
\langle K_a\rangle_{a\in A}\longmapsto\coprod_{a\in A}K_a=\langle L_{a,\alpha}\rangle_{(a,\alpha)\in\coprod_aA_a},
\]
exists and is left adjoint to the one-legged inclusion
$\iota_{\Fel(\cC)}\colon\Fel(\cC)\to\Fel^2(\cC)$. We call it \emph{flattening}.
It satisfies
\[
\body_{\Fel(\cC)}\circ\iota_{\Fel(\cC)}\iso\id_{\Fel(\cC)}\qquad\text{and}\qquad
\body_{\Fel(\cC)}\circ\Fel(\iota_\cC)\iso\id_{\Fel(\cC)}.
\]
Its restriction to $\Fel(\Kit(\cC))$ is the flattening of litters of
\cref{def:litter}.
\end{proposition}

\begin{proof}
The first assertions are \cref{prop:coproducts,prop:body-adjoint} applied to
$\Fel(\cC)$. The body of a unipedal cat of cats $\langle K\rangle_1$ is
$K$ reindexed along $1\times A\iso A$. For a cat $K=\langle L_\alpha\rangle_A$,
$\Fel(\iota_\cC)(K)=\langle\iota L_\alpha\rangle_A$, whose flattening is
$\langle L_\alpha\rangle$ indexed by $\coprod_\alpha1\iso A$. Both isomorphisms
have bijective gaits and identity components and are natural. The final
statement is a comparison of definitions.
\end{proof}

\begin{remark}\label{rem:pseudomonad}
\Cref{prop:flattening} supplies the unit $\iota$ and multiplication
$\body_{\Fel(\cdot)}$ of what is, on families, the well-known pseudomonad of
free coproduct completion, together with the two unit isomorphisms. We do not
verify the remaining coherence conditions of a pseudomonad, and we make no use
of them.
\end{remark}

\begin{theorem}\label{thm:iterated-point}
For every $n\geq1$ there is an equivalence $\Fel^n(\one)\eqv\Set^{[n-1]}$,
where $[n-1]=\{0<1<\dots<n-1\}$. Under it, an object of $\Fel^n(\one)$
corresponds to a chain of functions $X_0\to X_1\to\dots\to X_{n-1}$ in which
$X_{n-1}$ is the set of legs of the outermost cat.
\end{theorem}

\begin{proof}
For $n=1$ this is \cref{prop:Fel-point}, as $\Set^{[0]}=\Set$. Suppose
$\Fel^n(\one)\eqv\Set^{[n-1]}$. By \cref{cor:Fel-equivalence},
$\Fel^{n+1}(\one)\eqv\Fel(\Set^{[n-1]})$, and by \cref{thm:Fel-Set-J},
$\Fel(\Set^{[n-1]})\eqv\Set^{[n-1]^{\triangleright}}=\Set^{[n]}$, the adjoined
terminal object $\top=n$ carrying the leg set.
\end{proof}

\begin{theorem}\label{thm:iterated-presheaf}
Let $\cD$ be small, and let $\cD^{\triangleleft n}$ denote $\cD$ with $n$
successively adjoined initial objects, $\cD^{\triangleleft 0}=\cD$ and
$\cD^{\triangleleft(n+1)}=(\cD^{\triangleleft n})^{\triangleleft}$. Then
$\Fel^n(\PSh(\cD))\eqv\PSh(\cD^{\triangleleft n})$.
\end{theorem}

\begin{proof}
Induction on $n$, using \cref{cor:Fel-presheaf,cor:Fel-equivalence}.
\end{proof}

\begin{corollary}\label{cor:tower}
\begin{enumerate}
\item Let $\cD$ be another category with an initial and a terminal object. If
  $\Fel^n(\cC)\eqv\Fel^n(\cD)$ for some $n\geq0$, then $\cC\eqv\cD$.
\item For $m\neq n$, $\Fel^m(\one)\not\eqv\Fel^n(\one)$.
\end{enumerate}
\end{corollary}

\begin{proof}
(i) Induction on $n$, using \cref{thm:reconstruction} for the categories
$\Fel^{n-1}(\cC)$ and $\Fel^{n-1}(\cD)$, which have initial and terminal
objects. (ii) Suppose $m>n$ and $\Fel^m(\one)\eqv\Fel^n(\one)$. Since
$\Fel^m(\one)=\Fel^n(\Fel^{m-n}(\one))$, part (i) applied to
$\cC=\Fel^{m-n}(\one)$ and $\cD=\one$ gives $\Fel^{m-n}(\one)\eqv\one$. But for $k\geq1$ the category $\Fel^k(\one)$ has the
non-isomorphic objects $\Legless$ and $\Caudal$, so it is not equivalent to $\one$.
\end{proof}

\section{Gaited cats}\label{sec:gaits}

\begin{assumption}\label{ass:gaits}
In this section $\cC$ has an initial object $0$ and a terminal object $1$.
\end{assumption}

\begin{definition}\label{def:gaited-cat}
Let $\mathbb G$ be a small category, called a \emph{gait}. A
\emph{$\mathbb G$-gaited $\cC$-cat} is a tuple $(H,T,L)$ with $H$ initial, $T$
terminal and $L\colon\mathbb G\to\cC$ a functor, together with the unique
morphisms $H\to L(g)\to T$. A morphism $(H,T,L)\to(H',T',L')$ between cats with
gaits $\mathbb G$ and $\mathbb G'$ is a tuple $(u,F_H,F_T,\theta)$ with
$u\colon\mathbb G\to\mathbb G'$ a functor, $F_H\colon H\to H'$, $F_T\colon T\to T'$
and $\theta\colon L\Rightarrow L'\circ u$ a natural transformation, such that the
squares analogous to \eqref{eq:morphism-diagram} commute. We write
$\Fel_{\mathrm{gait}}(\cC)$ for the resulting category.
\end{definition}

A gaited cat is the same as a functor $(\mathbb G^{\triangleleft})^{\triangleright}\to\cC$
sending the adjoined initial and terminal objects to an initial and a terminal
object of $\cC$; for discrete $\mathbb G$ this is \cref{prop:cat-as-functor}.

Recall the category $\Diag(\cC)$ of small diagrams in $\cC$: objects are pairs
$(\mathbb G,L)$ with $\mathbb G$ small and $L\colon\mathbb G\to\cC$, and morphisms
$(\mathbb G,L)\to(\mathbb G',L')$ are pairs $(u,\theta)$ with $u\colon\mathbb G\to\mathbb G'$
and $\theta\colon L\Rightarrow L'u$, composed by $(v,\zeta)(u,\theta)=(vu,\zeta u\circ\theta)$.
It is the Grothendieck construction (\cref{def:groth-construction}) of the functor
$\Cat^{\op}\to\CAT$, $\mathbb G\mapsto\Fun(\mathbb G,\cC)$, $u\mapsto u^*$, and is
therefore fibred over $\Cat$.

\begin{proposition}\label{prop:gaited-diag}
The commutativity conditions in \cref{def:gaited-cat} are automatic, $F_H$ and
$F_T$ are unique isomorphisms, and the functor
$\Fel_{\mathrm{gait}}(\cC)\to\Diag(\cC)$ forgetting head and tail is an
equivalence. Its restriction to discrete gaits is the equivalence
$\Fel(\cC)\eqv\Fam(\cC)$ of \cref{thm:fundamental}, and $\Fam(\cC)$ is the full
subcategory of $\Diag(\cC)$ on discrete diagrams.
\end{proposition}

\begin{proof}
The proof of \cref{prop:automatic,thm:fundamental} applies verbatim. For the
last assertion, a functor between discrete categories is a function and a
natural transformation between diagrams on a discrete category is a family of
morphisms.
\end{proof}

\begin{proposition}\label{prop:gait-string}
Let $\legs_{\mathrm{gait}}\colon\Fel_{\mathrm{gait}}(\cC)\to\Cat$ send a gaited cat
to its gait. Let $\Delta_X\mathbb G$ be the gaited cat with gait $\mathbb G$ and
constant diagram at $X$, and let $\points_{\mathrm{gait}}(\mathbb G,L)$ be the
category of elements of the functor $\cC(1,L(-))\colon\mathbb G\to\Set$. Then
\[
\Delta_0\dashv\legs_{\mathrm{gait}}\dashv\Delta_1\dashv\points_{\mathrm{gait}}.
\]
If $\cC$ has small colimits, then the functor
$\colim\colon\Fel_{\mathrm{gait}}(\cC)\to\cC$ is left adjoint to the inclusion
$X\mapsto(\one,X)$ of unipedal cats, and it restricts on discrete gaits to the
body functor.
\end{proposition}

\begin{proof}
Using \cref{prop:gaited-diag} we compute in $\Diag(\cC)$. A morphism
$(\mathbb G,\Delta0)\to(\mathbb G',L')$ is a functor $u$ together with a natural
transformation $\Delta0\Rightarrow L'u$, of which there is exactly one; so
$\Diag(\cC)(\Delta_0\mathbb G,(\mathbb G',L'))\iso\Cat(\mathbb G,\mathbb G')$.
Dually $\Diag(\cC)((\mathbb G,L),\Delta_1\mathbb G')\iso\Cat(\mathbb G,\mathbb G')$.
A morphism $(\mathbb G,\Delta1)\to(\mathbb G',L')$ is a functor $u$ with a family
$x_g\in\cC(1,L'(ug))$ such that $L'(u(f))\circ x_g=x_{g'}$ for $f\colon g\to g'$;
this is precisely a functor $\mathbb G\to\points_{\mathrm{gait}}(\mathbb G',L')$,
$g\mapsto(ug,x_g)$, $f\mapsto u(f)$. These bijections are natural. Finally,
$\Diag(\cC)((\mathbb G,L),(\one,X))$ is the set of cocones from $L$ to $X$, which
is $\cC(\colim L,X)$.
\end{proof}

\begin{remark}\label{rem:gait-honest}
\Cref{prop:gait-string} specializes to \cref{thm:adjoint-string,prop:body-adjoint}
for discrete gaits, where the category of elements of $\cC(1,L(-))$ is the
discrete category on the set of points. All statements are standard facts
about the category of diagrams.
\end{remark}

\section{Two-dimensional cats}\label{sec:two-dim}

\begin{assumption}\label{ass:two-dim}
In this section $\cK$ is a strict $2$-category. We write $\cK(X,Y)$ for its
hom-categories, $\circ$ for composition of $1$-cells, $\cdot$ for vertical
composition of $2$-cells, and $*$ for horizontal composition.
\end{assumption}

\begin{definition}\label{def:bi-initial}
An object $H$ of $\cK$ is \emph{bi-initial} if every hom-category $\cK(H,X)$ is
equivalent to the terminal category $\one$, and \emph{bi-terminal} if every
$\cK(X,T)$ is equivalent to $\one$.
\end{definition}

\begin{lemma}\label{lem:contractible-hom}
A category is equivalent to $\one$ if and only if it is nonempty and each of its
hom-sets is a singleton. In particular, if $H$ is bi-initial, then for any two
$1$-cells $f,g\colon H\to X$ there is exactly one $2$-cell $f\Rightarrow g$,
and it is invertible.
\end{lemma}

\begin{proof}
A category $\cA$ is equivalent to $\one$ if and only if the unique functor
$\cA\to\one$ is an equivalence, that is, essentially surjective and fully
faithful. It is essentially surjective if and only if $\cA$ is nonempty, and
fully faithful if and only if every map $\cA(a,a')\to\one(*,*)$ is a bijection,
that is, every hom-set of $\cA$ is a singleton. The unique $2$-cell $f\Rightarrow g$ has the unique $2$-cell
$g\Rightarrow f$ as inverse.
\end{proof}

\begin{definition}\label{def:two-dim-cat}
A \emph{two-dimensional $\cK$-cat} is a tuple
\[
K=\bigl(H,T,(L_\alpha)_{\alpha\in A},(h_\alpha),(t_\alpha),\sigma\bigr)
\]
where $H$ is bi-initial, $T$ is bi-terminal, the $L_\alpha$ are objects, and
$h_\alpha\colon H\to L_\alpha$, $t_\alpha\colon L_\alpha\to T$ and
$\sigma\colon H\to T$ are chosen $1$-cells. By \cref{lem:contractible-hom}
there is a unique invertible $2$-cell $\omega_\alpha\colon t_\alpha\circ h_\alpha\Rightarrow\sigma$.
\end{definition}

Unlike the one-dimensional case, the canonical $1$-cells are now data: they
are unique only up to unique invertible $2$-cell.

\begin{definition}\label{def:two-dim-morphism}
A \emph{pseudo-morphism} $F\colon K\to K'$ of two-dimensional $\cK$-cats is a
tuple $(\varphi,F_H,F_T,(F_\alpha),(\lambda_\alpha),(\rho_\alpha),\kappa)$
consisting of a function $\varphi\colon A\to B$, $1$-cells $F_H\colon H\to H'$,
$F_T\colon T\to T'$, $F_\alpha\colon L_\alpha\to L'_{\varphi(\alpha)}$, and invertible
$2$-cells
\[
\lambda_\alpha\colon F_\alpha\circ h_\alpha\Rightarrow h'_{\varphi(\alpha)}\circ F_H,\qquad
\rho_\alpha\colon t'_{\varphi(\alpha)}\circ F_\alpha\Rightarrow F_T\circ t_\alpha,\qquad
\kappa\colon F_T\circ\sigma\Rightarrow\sigma'\circ F_H,
\]
such that, for each $\alpha$, the $2$-cell $t'_{\varphi(\alpha)}\circ F_\alpha\circ h_\alpha\Rightarrow\sigma'\circ F_H$
obtained by pasting $\lambda_\alpha$ and $\omega'_{\varphi(\alpha)}$ equals the one
obtained by pasting $\rho_\alpha$, $\omega_\alpha$ and $\kappa$. A
\emph{$2$-cell} $F\Rightarrow G$ between pseudo-morphisms with the \emph{same}
gait $\varphi$ is a tuple of $2$-cells $\Theta_H\colon F_H\Rightarrow G_H$,
$\Theta_T\colon F_T\Rightarrow G_T$, $\Theta_\alpha\colon F_\alpha\Rightarrow G_\alpha$
compatible with the $\lambda$, $\rho$ and $\kappa$; there are no $2$-cells
between pseudo-morphisms with different gaits. Composition is by composing
$1$-cells and pasting $2$-cells. We write $\Fel_2(\cK)$ for the resulting
$2$-category.
\end{definition}

\begin{proposition}[Automatic pseudo-commutativity]\label{prop:auto-pseudo}
In \cref{def:two-dim-morphism}:
\begin{enumerate}
\item for any choice of $\varphi$, $F_H$, $F_T$ and $(F_\alpha)$ there exist
  unique $2$-cells $\lambda_\alpha$, $\rho_\alpha$, $\kappa$, and they are
  invertible;
\item the coherence condition on a pseudo-morphism holds automatically;
\item for pseudo-morphisms $F,G$ with the same gait and any family of $2$-cells
  $\Theta_\alpha\colon F_\alpha\Rightarrow G_\alpha$, there are unique $\Theta_H$,
  $\Theta_T$, and the compatibility conditions hold automatically.
\end{enumerate}
\end{proposition}

\begin{proof}
The source and target $1$-cells of $\lambda_\alpha$ have the bi-initial domain
$H$; those of $\rho_\alpha$ have the bi-terminal codomain $T'$; those of $\kappa$
have domain $H$. By \cref{lem:contractible-hom} each of these $2$-cells exists,
is unique and is invertible. The coherence condition in (ii) is an equation
between two $2$-cells between the same pair of $1$-cells $H\to T'$, which holds
by uniqueness. In (iii), $\Theta_H$ is a $2$-cell between $1$-cells $H\to H'$
and $\Theta_T$ one between $1$-cells $T\to T'$, so both exist uniquely, and each
compatibility condition is an equation between $2$-cells whose underlying
$1$-cells have domain $H$ or codomain $T'$.
\end{proof}

Let $\Fam(\cK)$ be the $2$-category of small families in $\cK$: objects
$(A,(L_\alpha))$, $1$-cells $(\varphi,(F_\alpha))$, and $2$-cells
$(\varphi,F)\Rightarrow(\varphi,G)$ the families of $2$-cells
$\Theta_\alpha\colon F_\alpha\Rightarrow G_\alpha$ (none between $1$-cells with
different gaits).

\begin{theorem}\label{thm:two-dim}
The forgetful $2$-functor $U_2\colon\Fel_2(\cK)\to\Fam(\cK)$ is a local
equivalence: each functor of hom-categories
\[
\Fel_2(\cK)(K,K')\too\Fam(\cK)(U_2K,U_2K')
\]
is surjective on objects and fully faithful. It is surjective on objects if and only if $\cK$ has a bi-initial and
a bi-terminal object, in which case $U_2$ is a biequivalence.
\end{theorem}

\begin{proof}
Surjectivity on objects of the hom-functors: given $(\varphi,(F_\alpha))$, the
hom-categories $\cK(H,H')$ and $\cK(T,T')$ are nonempty, so $F_H$ and $F_T$
may be chosen, and \cref{prop:auto-pseudo}(i),(ii) supply the remaining data.
Full faithfulness is \cref{prop:auto-pseudo}(iii). For surjectivity on
objects, a two-dimensional cat over a given family requires a bi-initial and a
bi-terminal object, and conversely, given such objects, the $1$-cells
$h_\alpha$, $t_\alpha$, $\sigma$ can be chosen since the relevant hom-categories
are nonempty. A $2$-functor which is a local equivalence and surjective on
objects is a biequivalence \cite{JohnsonYau2021}.
\end{proof}

\begin{example}\label{ex:two-dim-cat}
Let $\cK=\Cat$, the $2$-category of small categories. The only bi-initial
object is the empty category: if $H$ is nonempty then $\Cat(H,\varnothing)$ is
empty. The bi-terminal objects are the categories equivalent to $\one$, for
instance the indiscrete category on any nonempty small set. Thus a
two-dimensional $\Cat$-cat may have as its tail the indiscrete category on
$\mathbb N$, whereas a one-dimensional $\Cat$-cat must have a tail with exactly
one object and one morphism. By \cref{thm:two-dim} this extra freedom is
invisible up to biequivalence.
\end{example}

\begin{remark}\label{rem:infinity}
One may similarly consider cats in an $(\infty,1)$-category, with the head
initial and the tail terminal in the $\infty$-categorical sense. The
commutativity data would then be parametrized by contractible spaces, and one
expects the $\infty$-category of such cats to be equivalent to the
$\infty$-category of set-indexed families. We have not written out a proof, and
we make no claim to this effect.
\end{remark}
